\section{Digital Control System Design in the $z$ Domain}\label{sec:design}
\lecdate{25 November and 2 December 2019}
% Figures of Chapter 5 (z-plane root locus, pole contours, proportional design).

% --- small response sketches -------------------------------------------------------
% \resp{x}{y}{type}: sketch of a response centred at (x,y); type = const, decay, grow,
% osc, damp, growosc (continuous) or the sampled kinds sdecay, sgrow, salt, saltdecay, saltgrow, sconst
\newcommand{\resp}[3]{%
\begin{scope}[shift={(#1,#2)},scale=0.55]
  \draw[gray,thin] (-0.1,0) -- (2.0,0);
  \draw[gray,thin] (0,-0.9) -- (0,0.9);
  \def\rtype{#3}
  \ifx\rtype\relax\fi
  \ifnum\pdfstrcmp{#3}{const}=0 \draw[itublue,thick] (0,0.5) -- (1.9,0.5);\fi
  \ifnum\pdfstrcmp{#3}{decay}=0 \draw[itublue,thick,domain=0:1.9,samples=25] plot (\x,{0.8*exp(-2.2*\x)});\fi
  \ifnum\pdfstrcmp{#3}{grow}=0 \draw[itublue,thick,domain=0:1.9,samples=25] plot (\x,{0.04*exp(2.0*\x)});\fi
  \ifnum\pdfstrcmp{#3}{osc}=0 \draw[itublue,thick,domain=0:1.9,samples=50] plot (\x,{0.6*sin(\x*540)});\fi
  \ifnum\pdfstrcmp{#3}{damp}=0 \draw[itublue,thick,domain=0:1.9,samples=60] plot (\x,{0.8*exp(-1.8*\x)*sin(\x*540)});\fi
  \ifnum\pdfstrcmp{#3}{growosc}=0 \draw[itublue,thick,domain=0:1.9,samples=60] plot (\x,{0.07*exp(1.6*\x)*sin(\x*540)});\fi
  \ifnum\pdfstrcmp{#3}{sdecay}=0 \foreach \k in {0,...,6}{\draw[itublue,thick] (\k*0.28,0) -- (\k*0.28,{0.8*exp(-0.45*\k)});\fill[itublue] (\k*0.28,{0.8*exp(-0.45*\k)}) circle (1.2pt);}\fi
  \ifnum\pdfstrcmp{#3}{sgrow}=0 \foreach \k in {0,...,6}{\draw[itublue,thick] (\k*0.28,0) -- (\k*0.28,{0.1*exp(0.5*\k)});\fill[itublue] (\k*0.28,{0.1*exp(0.5*\k)}) circle (1.2pt);}\fi
  \ifnum\pdfstrcmp{#3}{sconst}=0 \foreach \k in {0,...,6}{\draw[itublue,thick] (\k*0.28,0) -- (\k*0.28,0.5);\fill[itublue] (\k*0.28,0.5) circle (1.2pt);}\fi
  \ifnum\pdfstrcmp{#3}{salt}=0 \foreach \k in {0,...,6}{\pgfmathsetmacro{\v}{0.5*(-1)^\k}\draw[itublue,thick] (\k*0.28,0) -- (\k*0.28,\v);\fill[itublue] (\k*0.28,\v) circle (1.2pt);}\fi
  \ifnum\pdfstrcmp{#3}{saltdecay}=0 \foreach \k in {0,...,6}{\pgfmathsetmacro{\v}{0.8*exp(-0.45*\k)*(-1)^\k}\draw[itublue,thick] (\k*0.28,0) -- (\k*0.28,\v);\fill[itublue] (\k*0.28,\v) circle (1.2pt);}\fi
  \ifnum\pdfstrcmp{#3}{saltgrow}=0 \foreach \k in {0,...,6}{\pgfmathsetmacro{\v}{0.1*exp(0.5*\k)*(-1)^\k}\draw[itublue,thick] (\k*0.28,0) -- (\k*0.28,\v);\fill[itublue] (\k*0.28,\v) circle (1.2pt);}\fi
  \ifnum\pdfstrcmp{#3}{sdamp}=0 \foreach \k in {0,...,6}{\pgfmathsetmacro{\v}{0.8*exp(-0.35*\k)*sin(\k*80)}\draw[itublue,thick] (\k*0.28,0) -- (\k*0.28,\v);\fill[itublue] (\k*0.28,\v) circle (1.2pt);}\fi
  \ifnum\pdfstrcmp{#3}{sgrowosc}=0 \foreach \k in {0,...,6}{\pgfmathsetmacro{\v}{0.1*exp(0.45*\k)*sin(\k*80)}\draw[itublue,thick] (\k*0.28,0) -- (\k*0.28,\v);\fill[itublue] (\k*0.28,\v) circle (1.2pt);}\fi
\end{scope}}

% ---- p.63: s-plane pole locations and responses -----------------------------------
\newcommand{\figSpoles}{%
\begin{tikzpicture}[x=1cm,y=1cm]
  \draw[->] (-4.6,0) -- (4.6,0) node[right,lab]{$\sigma$};
  \draw[->] (0,-1.5) -- (0,3.4) node[above,lab]{$\jj\omega$};
  \node[pole,draw=boardgreen] at (-2.3,0) {}; \node[pole,draw=boardgreen] at (2.3,0) {};
  \node[pole,draw=boardgreen] at (0,0) {}; \node[pole,draw=boardgreen] at (0,2.2) {};
  \node[pole,draw=boardgreen] at (-3.6,2.3) {}; \node[pole,draw=boardgreen] at (-3.6,1.2) {};
  \node[pole,draw=boardgreen] at (2.4,2.0) {};
  \resp{-2.8}{-1.15}{decay}; \resp{-0.4}{-1.15}{const}; \resp{1.7}{-1.15}{grow};
  \resp{0.25}{2.1}{osc}; \resp{-4.45}{2.2}{damp}; \resp{-4.45}{1.1}{decay};
  \resp{3.0}{1.8}{growosc};
  \node[lab] at (-2.3,0.35) {$-a$};
\end{tikzpicture}}

% ---- p.63: z-plane pole locations and responses -----------------------------------
\newcommand{\figZpoles}{%
\begin{tikzpicture}[x=1cm,y=1cm]
  \draw[thick,itublue] (0,0) circle (2.0);
  \draw[->] (-5.3,0) -- (5.3,0) node[right,lab]{$\Real[z]$};
  \draw[->] (0,-2.4) -- (0,2.7) node[above,lab]{$\Imag[z]$};
  \foreach \px/\ty in {-4.4/saltgrow,-2.0/salt,-1.0/saltdecay,1.0/sdecay,2.0/sconst,3.4/sgrow}{
    \node[pole,draw=boardgreen] (p\ty) at (\px,0) {};
    \draw[->,gray] (\px,-0.15) -- (\px,-2.55);
    \resp{\px-0.5}{-3.45}{\ty};
  }
  \node[pole,draw=boardgreen] at (0.9,1.2) {}; \resp{1.0}{1.45}{sdamp};
  \node[pole,draw=boardgreen] at (-3.4,1.3) {}; \resp{-4.6}{1.55}{sgrowosc};
\end{tikzpicture}}

% ---- p.66: z = e^{sT}: |z| and angle; sigma const and wd const contours ----------
\newcommand{\figMapPoint}{%
\begin{tikzpicture}
  \draw[thick,gray] (0,0) circle (1.5);
  \draw[->] (-2.1,0) -- (2.3,0) node[right,lab]{$\Real[z]$};
  \draw[->] (0,-1.8) -- (0,2.0) node[above,lab]{$\Imag[z]$};
  \node[lab,below left] at (-1.5,0) {$-1$}; \node[lab,below right] at (1.5,0) {$1$};
  \coordinate (pz) at (0.82,1.02);
  \draw[boardgreen,thick] (0,0) -- (pz) node[midway,above left,lab]{$|z|$};
  \node[pole,draw=boardgreen] at (pz) {};
  \draw[boardgreen] (0.55,0) arc (0:51:0.55); \node[lab,boardgreen] at (0.85,0.28) {$\angle z$};
\end{tikzpicture}}

\newcommand{\figSigmaContours}{%
\begin{tikzpicture}
  \begin{scope}
    \draw[->] (-2.3,0) -- (2.3,0) node[right,lab]{$\sigma$};
    \draw[->] (0,-1.7) -- (0,1.8) node[above,lab]{$\jj\omega$};
    \foreach \sx in {-2.0,-1.5,-1.0,-0.5,0.5,1.0,1.5,2.0} \draw[gray] (\sx,-1.5) -- (\sx,1.5);
    \draw[gray] (0,-1.5) -- (0,1.5);
  \end{scope}
  \begin{scope}[xshift=6.2cm]
    \draw[->] (-2.3,0) -- (2.3,0) node[right,lab]{$\Real[z]$};
    \draw[->] (0,-1.9) -- (0,2.0) node[above,lab]{$\Imag[z]$};
    \foreach \rr in {0.15,0.4,0.7,1.0,1.35,1.7} \draw[boardgreen,thick] (0,0) circle (\rr);
  \end{scope}
\end{tikzpicture}}

\newcommand{\figWdContours}{%
\begin{tikzpicture}
  \begin{scope}
    \draw[->] (-2.4,0) -- (2.4,0) node[right,lab]{$\sigma$};
    \draw[->] (0,-2.3) -- (0,2.4) node[above,lab]{$\jj\omega$};
    \foreach \y/\l in {0.55/{\pi/4},1.1/{\pi/2},1.65/{3\pi/4}}{
      \draw[gray] (-2.2,\y) -- (2.2,\y) node[right,lab,boardgreen]{$\l$};
      \draw[gray] (-2.2,-\y) -- (2.2,-\y) node[right,lab,boardgreen]{$-\l$};}
  \end{scope}
  \begin{scope}[xshift=7.4cm]
    \draw[->] (-3.0,0) -- (3.1,0) node[right,lab]{$\Real[z]$};
    \draw[->] (0,-2.3) -- (0,2.9) node[above,lab]{$\Imag[z]$};
    \foreach \a/\l in {0/{0},45/{\pi/4},90/{\pi/2},135/{3\pi/4},180/{\pm\pi},-45/{-\pi/4},-90/{},-135/{-3\pi/4}}{
      \draw[boardgreen,thick] (0,0) -- (\a:2.1);
      \node[lab,boardgreen,font=\scriptsize] at (\a:2.4) {$\l$};}
  \end{scope}
\end{tikzpicture}}

% ---- p.65: second-order pole in the s plane ---------------------------------------
\newcommand{\figSecondOrderS}{%
\begin{tikzpicture}
  \draw[->] (-2.6,0) -- (1.8,0) node[right,lab]{$\sigma$};
  \draw[->] (0,-2.0) -- (0,2.1) node[above,lab]{$\jj\omega$};
  \coordinate (p) at (-1.5,1.4); \coordinate (q) at (-1.5,-1.4);
  \node[pole] at (p) {}; \node[pole] at (q) {};
  \draw[dashed,gray] (p) -- (-1.5,0) node[below left,lab]{$-\zeta\wn$};
  \draw[dashed,gray] (p) -- (0,1.4) node[right,lab]{$+\jj\wn\sqrt{1-\zeta^{2}}=\jj\wdd$};
  \draw[dashed,gray] (q) -- (-1.5,0); \draw[dashed,gray] (q) -- (0,-1.4) node[right,lab]{$-\jj\wdd$};
  \draw[boardgreen,thick] (0,0) -- (p) node[midway,above right,lab]{$\wn$};
  \draw[boardgreen] (-0.5,0) arc (180:137:0.5); \node[lab,boardgreen] at (-0.75,0.28) {$\theta$};
\end{tikzpicture}}

% ---- z-plane grid: zeta spirals and wn contours (zgrid) -----------------------------
% the same curves as MATLAB's zgrid; drawn for the upper half plane and mirrored
\newcommand{\figZgrid}{%
\begin{tikzpicture}[scale=3.1]
  \clip (-1.18,-1.12) rectangle (1.3,1.12);
  \draw[gray] (0,0) circle (1);
  \draw[->] (-1.15,0) -- (1.28,0) node[right,lab]{$\Real[z]$};
  \draw[->] (0,-1.1) -- (0,1.1) node[above,lab]{$\Imag[z]$};
  % zeta spirals: |z| = exp(-zeta*th/sqrt(1-zeta^2)), th = wd T in [0,pi]
  \foreach \zt in {0.1,0.2,0.3,0.4,0.5,0.6,0.7,0.8,0.9}{
    \draw[boardgreen,thick] plot[domain=0:180,samples=60,variable=\th]
      ({exp(-\zt*\th*pi/180/sqrt(1-\zt*\zt))*cos(\th)},{exp(-\zt*\th*pi/180/sqrt(1-\zt*\zt))*sin(\th)});
    \draw[boardgreen,thick] plot[domain=0:180,samples=60,variable=\th]
      ({exp(-\zt*\th*pi/180/sqrt(1-\zt*\zt))*cos(\th)},{-exp(-\zt*\th*pi/180/sqrt(1-\zt*\zt))*sin(\th)});
  }
  % wn contours: |z| = exp(-sqrt((wnT)^2 - th^2)), th from 0 to wnT (rad), wnT = k pi/10
  \foreach \kk in {1,2,...,9}{
    \pgfmathsetmacro{\wnT}{\kk*pi/10}
    \draw[gray,thick] plot[domain=0:\wnT*180/pi,samples=40,variable=\th]
      ({exp(-sqrt(max(0,\wnT*\wnT-(\th*pi/180)^2)))*cos(\th)},{exp(-sqrt(max(0,\wnT*\wnT-(\th*pi/180)^2)))*sin(\th)});
    \draw[gray,thick] plot[domain=0:\wnT*180/pi,samples=40,variable=\th]
      ({exp(-sqrt(max(0,\wnT*\wnT-(\th*pi/180)^2)))*cos(\th)},{-exp(-sqrt(max(0,\wnT*\wnT-(\th*pi/180)^2)))*sin(\th)});
  }
  \node[lab,boardgreen,fill=white,inner sep=1pt] at (-0.78,0.60) {$\zeta=0.1$};
  \node[lab,boardgreen,fill=white,inner sep=1pt] at (0.40,0.15) {$0.9$};
  \node[lab,fill=white,inner sep=1pt,font=\scriptsize] at (0.78,0.93) {$0.3\pi/T$};
  \node[lab,fill=white,inner sep=1pt,font=\scriptsize] at (0.98,0.58) {$0.2\pi/T$};
  \node[lab,fill=white,inner sep=1pt,font=\scriptsize] at (1.12,0.30) {$0.1\pi/T$};
\end{tikzpicture}}

% ---- p.62: root loci -------------------------------------------------------------
\newcommand{\figRLfirst}{%
\begin{tikzpicture}
  \begin{axis}[width=0.52\linewidth,height=5.2cm,axis lines=middle,xlabel={$\Real[z]$},ylabel={$\Imag[z]$},
      xmin=-2.6,xmax=2.2,ymin=-1.7,ymax=1.7,axis equal image,xtick={-1,1},ytick=\empty,
      every axis x label/.style={at={(ticklabel* cs:1.0)},anchor=west},
      every axis y label/.style={at={(ticklabel* cs:1.0)},anchor=south},clip=false]
    \draw[gray,thick] (axis cs:0,0) circle[radius=1];
    \draw[thick,itublue] (axis cs:1,0) -- (axis cs:-2.4,0);
    \addplot[only marks,mark=x,mark size=3.5pt,very thick,itublue] coordinates {(1,0)};
    \addplot[only marks,mark=*,mark size=1.8pt,boardred] coordinates {(-1,0)};
    \node[lab,boardred,below left] at (axis cs:-1,-0.05) {$K_{cr}=2$};
  \end{axis}
\end{tikzpicture}}

\newcommand{\figRLsecond}{%
\begin{tikzpicture}
  \begin{axis}[width=0.78\linewidth,height=6.4cm,axis lines=middle,xlabel={$\Real[z]$},ylabel={$\Imag[z]$},
      xmin=-1.35,xmax=1.5,ymin=-1.25,ymax=1.25,axis equal image,xtick={-1,0.5,1},xticklabel style={font=\scriptsize},ytick=\empty,
      every axis x label/.style={at={(ticklabel* cs:1.0)},anchor=west},
      every axis y label/.style={at={(ticklabel* cs:1.0)},anchor=south},clip=false]
    \draw[gray,thick] (axis cs:0,0) circle[radius=1];
    % real branches: z^2-1.5z+0.5+K=0 for 0<K<0.0625 -> z = 0.75 +- sqrt(0.0625-K)
    \addplot[thick,itublue,domain=0:0.0625,samples=30] ({0.75+sqrt(0.0625-x)},0);
    \addplot[thick,itublue,domain=0:0.0625,samples=30] ({0.75-sqrt(0.0625-x)},0);
    % complex branches for K>0.0625 : z=0.75 +- j sqrt(K-0.0625) (vertical line)
    \addplot[thick,itublue,domain=0:1.2,samples=2] (0.75,x);
    \addplot[thick,itublue,domain=0:1.2,samples=2] (0.75,-x);
    \addplot[only marks,mark=x,mark size=3.5pt,very thick,itublue] coordinates {(0.5,0)(1,0)};
    \addplot[only marks,mark=*,mark size=1.8pt,boardred] coordinates {(0.75,0.6614)(0.75,-0.6614)};
    \node[lab,boardred,right] at (axis cs:0.78,0.74) {$K=0.5$};
    \draw[dashed,gray] (axis cs:0.75,-1.2) -- (axis cs:0.75,1.2);
    \node[lab,gray,below] at (axis cs:0.75,-1.2) {$\sigma_a=0.75$};
  \end{axis}
\end{tikzpicture}}

\newcommand{\zspiral}[1]{\addplot[gray,dashed,domain=0:80,samples=40]
  ({exp(-#1*rad(x)/sqrt(1-#1*#1))*cos(x)},{exp(-#1*rad(x)/sqrt(1-#1*#1))*sin(x)});}

% ---- design example: root locus with design points -----------------------------------
\newcommand{\figRLdesign}{%
\begin{tikzpicture}
  \begin{axis}[width=\linewidth,height=7.4cm,axis lines=middle,xlabel={$\Real[z]$},ylabel={$\Imag[z]$},
      xmin=-0.3,xmax=1.15,ymin=-0.62,ymax=0.62,axis equal image,xtick={0,0.5,1},ytick={-0.5,0.5},
      xticklabel style={font=\scriptsize},yticklabel style={font=\scriptsize},clip=false,
      every axis x label/.style={at={(ticklabel* cs:1.0)},anchor=west},
      every axis y label/.style={at={(ticklabel* cs:1.0)},anchor=south}]
    \addplot[gray,thick,domain=-0.3:1.15,samples=2,draw=none]{0};
    \addplot[gray,thick,domain=0:40,samples=40]({cos(x)},{sin(x)});
    % zeta contours 0.3, 0.436, 0.7 (spirals)
    \zspiral{0.3}\zspiral{0.436}\zspiral{0.7}
    % constant radius |z|=0.82 circle (tau = 0.5 s)
    \addplot[boardgreen,dashed,domain=0:90,samples=40]({0.8187*cos(x)},{0.8187*sin(x)});
    \addplot[thick,itublue,domain=0:0.0625,samples=30] ({0.75+sqrt(0.0625-x)},0);
    \addplot[thick,itublue,domain=0:0.0625,samples=30] ({0.75-sqrt(0.0625-x)},0);
    \addplot[thick,itublue,domain=0:0.62,samples=2] (0.75,x);
    \addplot[thick,itublue,domain=0:0.62,samples=2] (0.75,-x);
    \addplot[only marks,mark=x,mark size=3.5pt,very thick,itublue] coordinates {(0.5,0)(1,0)};
    \addplot[only marks,mark=*,mark size=1.8pt,boardred] coordinates {(0.75,0.4097)(0.75,0.3279)(0.75,0.1936)};
    \node[lab,boardred,left,fill=white,inner sep=1pt] at (axis cs:0.73,0.44) {a) $K=0.23$};
    \node[lab,boardred,left,fill=white,inner sep=1pt] at (axis cs:0.73,0.33) {b) $K=0.17$};
    \node[lab,boardred,left,fill=white,inner sep=1pt] at (axis cs:0.73,0.19) {c) $K=0.10$};
  \end{axis}
\end{tikzpicture}}


There are two ways to design a digital controller:
\begin{enumerate}
  \item The digital controller can be designed directly in the $z$ domain (this chapter).
  \item The digital controller can be obtained from the model of an analog controller (next chapter).
\end{enumerate}

%==============================================================================
\subsection{$z$-domain root locus}\scanref{pp.\,61--63}
The root locus method provides information about the closed-loop system behaviour based on its open-loop
poles and zeros. In continuous time,
\[
  1+K\,L(s)=0,\qquad K:\ \text{design parameter},\quad L(s):\ \text{loop gain.}
\]
Two conditions and the related rules are used to obtain root-locus plots:
\begin{align*}
  \text{Magnitude condition:}&\quad K\,|L(s)|=1,\\
  \text{Angle condition:}&\quad \angle L(s)=\pm(2m+1)\,180^{\circ},\qquad m=0,1,2,\dots
\end{align*}
In discrete time,
\begin{equation}
  1+K\,L(z)=0,\qquad L(z)=C(z)\,\GZAS(z),\qquad C(z)=K\,\frac{N_c(z)}{D_c(z)}.
\end{equation}
All the rules derived for continuous time are applicable to discrete-time systems:
\begin{itemize}
  \item The number of root-locus branches equals the number of open-loop poles of $L(z)$.
  \item Root-locus branches start at open-loop poles and end at open-loop zeros or at $\infty$.
  \item The real-axis root loci have an odd number of poles plus zeros to their right (a point of the real
        axis belongs to the locus if the number of real poles and zeros to its right is odd).
\end{itemize}
\begin{itemize}
  \item The branches going to infinity asymptotically approach straight lines defined by the angle
  \begin{equation}
    \phi_a=\pm\frac{(2m+1)\,180^{\circ}}{n_p-n_z},\qquad m=0,1,2,\dots
  \end{equation}
  and the intercept at the point
  \begin{equation}
    \sigma_a=\frac{\sum p-\sum z}{n_p-n_z}.
  \end{equation}
  \item Breakaway points (points of departure from the real axis) correspond to local maxima of $K$, whereas
  break-in points (points of arrival at the real axis) correspond to local minima of $K$.
  \item The angle of departure from a complex pole $p_n$ is given by
  \[
    180^{\circ}-\sum(\angle p_n-\angle p_i)+\sum(\angle p_n-\angle z_i).
  \]
\end{itemize}

\begin{example}[First-order type-1 system]
Obtain the root locus plot and the critical gain $K_{cr}$ for the first-order type-1 system with the loop gain
$L(z)=\dfrac{1}{z-1}$.

\begin{figure}[H]
\centering
\figRLfirst
\caption{Root locus of $L(z)=1/(z-1)$ with the critical gain at $z=-1$.}
\end{figure}
\noindent The root locus lies entirely on the real axis between the open-loop pole $z_1=1$ and $-\infty$
(for $K>0$ the closed-loop pole is $z=1-K$). At the stability limit $z_1=-1$, the magnitude condition gives
\[
  K\,|L(z)|=1\ \to\ K\left|\frac{1}{z-1}\right|_{z=-1}=1\ \to\ K_{cr}=2.
\]
\end{example}

\begin{boardnote}
The board states that the locus lies on the real axis ``between open-loop pole and $\infty$''. It is the part
to the left of the pole, i.e.\ towards $-\infty$, and it crosses the unit circle at $z=-1$. This is the one
the sketch shows.
\end{boardnote}

\begin{example}[Second-order type-1 system]
Same question with $L(z)=\dfrac{1}{(z-1)(z-0.5)}$.

The open-loop poles are $z_1=1$ and $z_2=0.5$, with no finite zero ($n_p=2$, $n_z=0$):
\[
  \phi_a=\pm\frac{(2m+1)\,180^{\circ}}{2-0}=\pm90^{\circ},\qquad
  \sigma_a=\frac{1+0.5}{2-0}=0.75.
\]
The characteristic equation is
\[
  1+K\,L(z)=0\ \to\ 1+\frac{K}{(z-1)(z-0.5)}=0\ \to\ (z-1)(z-0.5)+K=z^{2}-1.5z+K+0.5=0.
\]
The locus leaves the real axis at the breakaway point $z=0.75$ (where $K=-(z-1)(z-0.5)$ is at a maximum,
$K=0.0625$) and then follows the vertical asymptote $\Real[z]=0.75$. It crosses the unit circle where
$|z_{1,2}|^{2}=1$. For the complex pair $z_{1,2}=a\pm\jj b$ the characteristic polynomial is
$z^{2}-2a\,z+a^{2}+b^{2}$, and $|z_{1,2}|^{2}=\bigl(\sqrt{a^{2}+b^{2}}\bigr)^{2}=a^{2}+b^{2}$, hence
\[
  |z_{1,2}|^{2}=1=K+0.5\ \Rightarrow\ K=0.5,\qquad z_{1,2}=0.75\pm\jj0.661.
\]
\begin{figure}[H]
\centering
\figRLsecond
\caption{Root locus of $L(z)=1/\bigl((z-1)(z-0.5)\bigr)$; the locus crosses the unit circle at $K=0.5$.}
\end{figure}
\end{example}

\begin{boardnote}
The board writes the asymptote angles as $\pm90^{\circ},\ \pm180^{\circ}$ and the intercept as
$\dfrac{-1-0.5}{2-0}=-0.75$. With the open-loop poles at $+1$ and $+0.5$, two asymptotes exist
($n_p-n_z=2$), at $\pm90^{\circ}$, and $\sigma_a=\dfrac{1+0.5}{2}=+0.75$; this is also where the board draws
the vertical line in its figure. The typeset values above are therefore $\pm90^{\circ}$ and $+0.75$.
\end{boardnote}

\addedsection[Added problems: root locus in the $z$ plane]
\begin{problem}[Critical gain]
For $L(z)=\dfrac{1}{(z-1)(z-0.5)}$ verify $K_{cr}=0.5$ with the Jury conditions, and find the gain at the
breakaway point.
\end{problem}
\begin{solution}
The characteristic polynomial is $z^{2}-1.5z+(0.5+K)$ with $a_2=1$, $a_1=-1.5$, $a_0=0.5+K$:
$a_2-a_1+a_0=3+K>0$, $a_2+a_1+a_0=K>0$, $|a_0|<a_2\Rightarrow0.5+K<1\Rightarrow K<0.5$. So stable for
$0<K<0.5$: $K_{cr}=0.5$, matching the locus. The breakaway point is at $z=0.75$ (the maximum of
$K(z)=-(z-1)(z-0.5)$), with $K=0.25\cdot0.25=0.0625$: for $K<0.0625$ the closed-loop poles are real, for larger
$K$ complex.
\end{solution}

%------------------------------------------------------------------------------
\subsection{$z$-domain pole locations and time responses}\scanref{pp.\,63--64}
\begin{figure}[H]
\centering
\resizebox{0.95\linewidth}{!}{\figSpoles}\\[2mm]
\resizebox{0.95\linewidth}{!}{\figZpoles}
\caption{Pole locations in the $s$ domain and in the $z$ domain, and the associated time responses
(sketches).}
\end{figure}

\begin{itemize}
  \item Positive real poles are associated with exponentials, like the continuous case.
  \item Negative real poles are associated with sequences of alternating signals.
\end{itemize}
To obtain the desired relationship between $s$-domain and $z$-domain poles we examine two key cases (a real pole
and a complex-conjugate pair of poles). More complex time functions can be reduced to these cases by
partial-fraction expansion.

\begin{center}
\resizebox{\linewidth}{!}{%
\renewcommand{\arraystretch}{2.2}
\begin{tabular}{@{}c c c c@{}}
Continuous & Laplace & Sampled & $z$ transform\\\hline
$f(t)=\ee^{-\alpha t},\ t\ge0$ & $F(s)=\dfrac{1}{s+\alpha}$ & $f(kT)=\ee^{-\alpha kT},\ k\ge0$ & $F(z)=\dfrac{z}{z-\ee^{-\alpha T}}$\\
$f(t)=\ee^{-\alpha t}\sin(\wdd t)$ & $F(s)=\dfrac{\wdd}{(s+\alpha)^{2}+\wdd^{2}}$ & $f(kT)=\ee^{-\alpha kT}\sin(\wdd kT)$ &
$F(z)=\dfrac{\sin(\wdd T)\,\ee^{-\alpha T}\,z}{z^{2}-2\cos(\wdd T)\,\ee^{-\alpha T}z+\ee^{-2\alpha T}}$
\end{tabular}}
\end{center}
Using the table, it appears that if $F(s)$ has a pole at $-\alpha$, then $F(z)$ has a pole at $\ee^{-\alpha T}$;
if $F(s)$ has a pole at $(-\zeta\wn+\jj\wdd)$, then $F(z)$ has a pole at $\ee^{(-\zeta\wn+\jj\wdd)T}$.
We conclude that
\begin{equation}
  p_z=\ee^{p_sT}\qquad(\text{like }z=\ee^{sT}).
\end{equation}
If $p_s=\sigma\pm\jj\wdd$, then $p_z=\ee^{(\sigma\pm\jj\wdd)T}=\ee^{\sigma T}\,\ee^{\pm\jj\wdd T}=|z|\,\ee^{\pm\jj\angle z}$, so
\begin{equation}
  |z|=\ee^{\sigma T},\qquad \angle z=\wdd T.
\end{equation}

\addedsection[Added problem: mapping $s$ poles to $z$ poles]
\begin{problem}
Map the $s$-plane pole $s=-2+\jj4$ into the $z$ plane for $T=0.1\,\mathrm{s}$ and $T=0.5\,\mathrm{s}$, and
comment.
\end{problem}
\begin{solution}
$z=\ee^{sT}$: for $T=0.1$, $|z|=\ee^{-0.2}=0.8187$ and $\angle z=0.4\,\mathrm{rad}=22.9^{\circ}$, i.e.\
$z=0.754+0.319\,\jj$. For $T=0.5$, $|z|=\ee^{-1}=0.368$ and $\angle z=2\,\mathrm{rad}=114.6^{\circ}$, i.e.\
$z=-0.153+0.335\,\jj$. The same analog pole produces a very different $z$ pole as $T$ grows: the angle
$\wdd T$ increases (fewer samples per oscillation) and the radius shrinks.
\end{solution}

%------------------------------------------------------------------------------
\subsection{$z$-domain contours}\scanref{pp.\,65--67}
\paragraph{Reminder: the second-order system in continuous time.} The characteristic equation of a general
second-order system is
\[
  s^{2}+2\zeta\wn s+\wn^{2}=s^{2}+2\zeta\wn s+\zeta^{2}\wn^{2}+\wn^{2}(1-\zeta^{2})
  =(s+\zeta\wn)^{2}+\underbrace{\wn^{2}(1-\zeta^{2})}_{\wdd^{2}},
\]
\[
  \Rightarrow\ s_{1,2}=-\zeta\wn\pm\jj\wn\sqrt{1-\zeta^{2}}=-\zeta\wn\pm\jj\wdd.
\]
\begin{figure}[H]
\centering
\begin{minipage}{0.5\linewidth}\figSecondOrderS\end{minipage}\hfill
\begin{minipage}{0.45\linewidth}
$\wn$: natural frequency,\\ $\wdd$: damped natural frequency,\\ $\zeta$: damping ratio, $\zeta=\cos\theta$.
\end{minipage}
\caption{Poles of the second-order system in the $s$ plane.}
\end{figure}

\paragraph{The $z$ plane contours.} With $s=\sigma+\jj\wdd$ and $z=\ee^{sT}$,
\begin{equation}
  z=\ee^{(\sigma\pm\jj\wdd)T}=\ee^{\sigma T}\,\ee^{\pm\jj\wdd T}=|z|\,\ee^{\pm\jj\angle z},\qquad
  |z|=\ee^{\sigma T}\ \text{(magnitude)},\quad \angle z=\wdd T\ \text{(angle)}.
\end{equation}
\begin{figure}[H]
\centering
\resizebox{0.35\linewidth}{!}{\figMapPoint}
\caption{A pole in the $z$ plane: magnitude $|z|=\ee^{\sigma T}$ and angle $\angle z=\wdd T$.}
\end{figure}

\begin{center}
\renewcommand{\arraystretch}{1.5}
\begin{tabular}{@{}l l l l@{}}
Contour & $s$-domain poles & Contour & $z$-domain poles\\\hline
$\sigma=$ constant & vertical lines & $|z|=\ee^{\sigma T}=$ const. & circles\\
$\wdd=$ constant & horizontal lines & $\angle z=\wdd T=$ const. & radial lines
\end{tabular}
\end{center}

\begin{figure}[H]
\centering
\resizebox{0.92\linewidth}{!}{\figSigmaContours}
\caption{Constant-$\sigma$ contours in the $s$ plane and in the $z$ plane.}
\end{figure}
\begin{figure}[H]
\centering
\resizebox{0.92\linewidth}{!}{\figWdContours}
\caption{Constant-$\wdd$ contours in the $s$ plane and in the $z$ plane (the labelled values are the angles
$\wdd T$ in radians, here for $\wdd T=\pm\pi/4,\pm\pi/2,\pm3\pi/4$).}
\end{figure}

\paragraph{Remark.} In the $s$ domain the $\zeta$ lines and the $\wn$ contours are also very important for
designing control systems. For a root-locus MATLAB plot these lines can be added by the \texttt{sgrid} command.
\begin{center}
\renewcommand{\arraystretch}{1.3}
\begin{tabular}{@{}l l l@{}}
$\zeta>1$ & real and unequal & overdamped\\
$\zeta=1$ & real and equal & critically damped\\
$0<\zeta<1$ & complex conjugate & underdamped\\
$\zeta=0$ & purely imaginary & undamped
\end{tabular}
\end{center}

%------------------------------------------------------------------------------
\subsection{$\zeta$ and $\wn$ contours in the $z$ domain}\scanref{p.\,68}
The $z$-domain poles for a second-order underdamped system are
\[
  z_{1,2}=\ee^{(-\zeta\wn\pm\jj\wdd)T}.
\]
Let $\theta=\wdd T=\wn T\sqrt{1-\zeta^{2}}$, so that $\wn T=\dfrac{\theta}{\sqrt{1-\zeta^{2}}}$. Then
\[
  |z|=\ee^{-\zeta\wn T}=\ee^{-\zeta\,\theta/\sqrt{1-\zeta^{2}}}\qquad(\theta\ \text{in radians}),
\]
which is the \emph{spiral equation}: the constant-$\zeta$ contours are logarithmic spirals. With $\theta$ in
degrees, as on the board,
\begin{equation}
  |z|=\exp\!\left(-\frac{\zeta\,(\pi/180^{\circ})\,\theta}{\sqrt{1-\zeta^{2}}}\right).
\end{equation}
For constant $\wn$: from $\theta=\wn T\sqrt{1-\zeta^{2}}$ we have $\zeta^{2}=1-\left(\dfrac{\theta}{\wn T}\right)^{2}$ and
\begin{equation}
  |z|=\ee^{-\zeta\wn T}=\ee^{-\sqrt{(\wn T)^{2}-\theta^{2}}}.
\end{equation}
The \texttt{zgrid} command of MATLAB can be used to obtain the $\zeta$ and $\wn$ contours.

\begin{figure}[H]
\centering
\figZgrid
\caption{The $\zeta$ contours (green spirals, $\zeta=0.1,\dots,0.9$) and the $\wn$ contours (grey,
$\wn=0.1\pi/T,\,0.2\pi/T,\dots,0.9\pi/T$) in the $z$ plane, as produced by \texttt{zgrid}.}
\end{figure}

\addedsection[Added problem: reading $\zeta$ and $\wn$ from a pole]
\begin{problem}
A closed-loop pole pair is at $z=0.75\pm\jj0.41$ with $T=0.1\,\mathrm{s}$. Find $\zeta$, $\wn$ and the
settling time.
\end{problem}
\begin{solution}
$|z|=\sqrt{0.75^{2}+0.41^{2}}=0.8548$ and $\theta=\wdd T=\arctan(0.41/0.75)=0.5003\,\mathrm{rad}$, so
$\wdd=5.0\,\mathrm{rad/s}$ and $\zeta\wn=-\ln|z|/T=1.573$. Then $\wn=\sqrt{\wdd^{2}+(\zeta\wn)^{2}}=5.24\,\mathrm{rad/s}$,
$\zeta=1.573/5.24=0.30$ and $T_s=4/(\zeta\wn)=2.54\,\mathrm{s}$. These are the values of design a) below.
\end{solution}

%------------------------------------------------------------------------------
\subsection{Design specifications and proportional control design}\scanref{p.\,69}
The specifications for $z$-domain design are similar to those for $s$-domain design, such as
\begin{align*}
  \text{Time constant}&\ \to\ \tau=\frac{1}{\zeta\wn},\\
  \text{Settling time (2\%)}&\ \to\ T_s=\frac{4}{\zeta\wn},\\
  \text{Percentage overshoot}&\ \to\ PO=\ee^{-\pi\zeta/\sqrt{1-\zeta^{2}}}\times100,\\
  \text{Peak time}&\ \to\ T_p=\frac{\pi}{\wdd}.
\end{align*}

\paragraph{Proportional control design in the $z$ domain.}
\begin{itemize}
  \item Proportional control design involves the selection of a DC gain to satisfy a time-response design
        specification.
  \item As in $s$-domain design, it is obtained by tuning the gain to select a dominant closed-loop pair in the
        appropriate region of the complex plane.
  \item Analytical design is possible for low-order systems but it is more difficult than its analog counterpart.
\end{itemize}

\begin{example}[Proportional controller design]
Design a proportional controller for the digital system $L(z)=\dfrac{1}{(z-1)(z-0.5)}$ with $T=0.1\,\mathrm{s}$
to obtain
\begin{enumerate}
  \item[a)] a damped natural frequency $\wdd$ of $5\,\mathrm{rad/s}$,
  \item[b)] a time constant $\tau$ of $0.5\,\mathrm{s}$,
  \item[c)] a damping ratio of $0.7$.
\end{enumerate}

\paragraph{Graphical approach.}
\begin{enumerate}
  \item[a)] $\wdd=5\,\mathrm{rad/s}$: the angle of the pole is $\wdd T=5\cdot0.1=0.5\,\mathrm{rad}$ (a radial line).
  \item[b)] $\tau=\dfrac{1}{\zeta\wn}=0.5\ \Rightarrow\ \zeta\wn=2$. With $z=\ee^{-(\zeta\wn\pm\jj\wdd)T}$,
  \[
    |z|=\ee^{-\zeta\wn T}=\ee^{-2(0.1)}=0.82,
  \]
  which shows a circle with radius $\ee^{-\zeta\wn T}$ (parametrised by $x=r\cos\theta$, $y=r\sin\theta$,
  $\theta=0:2\pi$).
  \item[c)] $\zeta=0.7$: we have directly the spirals on the $z$ plane to consider.
\end{enumerate}
By using the intersection of the lines and curves in the root locus plot we can calculate each $K$ value:
\begin{center}
\renewcommand{\arraystretch}{1.3}
\begin{tabular}{c c c c}
Design & Gain $K$ & $\zeta$ & $\wn$ [rad/s]\\\hline
a) & 0.23 & 0.3 & 5.24\\
b) & 0.17 & 0.4 & 4.60\\
c) & 0.10 & 0.7 & 3.63
\end{tabular}
\end{center}

\begin{figure}[H]
\centering
\figRLdesign
\caption{Root locus with the three design points; the dashed curves are the $\zeta$ spirals for
$0.3,\,0.436,\,0.7$ and the circle $|z|=0.82$.}
\end{figure}

\paragraph{Analytical approach.} The characteristic equation is
\[
  1+K\,L(z)=0\ \to\ z^{2}-1.5z+K+0.5=0
\]
and it must match the standard second-order equation
\[
  z^{2}-2\cos(\wdd T)\,\ee^{-\zeta\wn T}\,z+\ee^{-2\zeta\wn T}.
\]
We have two equations:
\[
  z^{1}:\ 1.5=2\cos(\wdd T)\,\ee^{-\zeta\wn T},\qquad z^{0}:\ K+0.5=\ee^{-2\zeta\wn T}.
\]
\begin{enumerate}
  \item[a)] $\wdd T=5\cdot0.1=0.5\,\mathrm{rad}$:
  \[
    1.5=2\cos(0.5)\,\ee^{-\zeta\wn(0.1)}\ \Rightarrow\ \zeta\wn=1.571\quad(*),\qquad
    K=\ee^{-2(1.571)(0.1)}-0.5=0.23.
  \]
  Using $\wdd T=\wn T\sqrt{1-\zeta^{2}}$: $0.5=0.1\,\wn\sqrt{1-\zeta^{2}}$ $(**)$. By $(*)$ and $(**)$,
  $\zeta=0.3$ and $\wn=5.24\,\mathrm{rad/s}$.
  \item[b)] $\tau=\dfrac{1}{\zeta\wn}=0.5\ \to\ \zeta\wn=2$ $(*)$:
  \[
    K+0.5=\ee^{-2(2)(0.1)}\ \Rightarrow\ K=0.17.
  \]
  \[
    1.5=2\cos(\wdd T)\ee^{-\zeta\wn T}\ \Rightarrow\ \wdd T=\cos^{-1}\!\left(\frac{1.5}{2\ee^{-\zeta\wn T}}\right)=0.4127,
  \]
  and $\wdd T=\wn T\sqrt{1-\zeta^{2}}$, i.e.\ $0.4127=\wn\cdot0.1\sqrt{1-\zeta^{2}}$ $(**)$. By $(*)$ and $(**)$:
  $\zeta=0.436$ and $\wn=4.586\,\mathrm{rad/s}$.
  \item[c)] $\zeta=0.7$: $1.5=2\cos\bigl(\wn T\sqrt{1-\zeta^{2}}\bigr)\ee^{-\zeta\wn T}$ gives $\wn=3.63\,\mathrm{rad/s}$ and
  $K+0.5=\ee^{-2\zeta\wn T}$ gives $K=0.10$.
\end{enumerate}
\end{example}

\begin{boardnote}
(1) In part a) the board's first line (page~70) writes $\wdd=0.5\,\mathrm{rad/s}$ and $\wdd T=0.5\cdot0.1=0.05\,\mathrm{rad}$,
while the working on the next page (and the answer $\wn=5.24\,\mathrm{rad/s}$) uses $\wdd=5\,\mathrm{rad/s}$ and
$\wdd T=0.5\,\mathrm{rad}$, as in the problem statement. (2) On page~63 the characteristic equation is
$z^{2}-1.5z+K+0.5$; on page~71 the second equation is written ``$K-0.5=\ee^{-2\zeta\wn T}$'', but the
numerical working ($K=\ee^{-2(1.571)(0.1)}-0.5=0.23$) is the one for $K+0.5$. (3) In part b) the table gives $\wn=4.60$ while the working on page~71 gives $4.586$; the computed value is
$4.586$. The remaining numbers agree with a direct numerical solution ($K=0.230,\ 0.170,\ 0.102$).
\end{boardnote}

\addedsection[Added problems: proportional design]
\begin{problem}[Design for a given damping ratio]
For the same plant ($L=1/((z-1)(z-0.5))$, $T=0.1\,\mathrm{s}$) find the proportional gain that gives
$\zeta=0.5$ and determine $\wn$, $\wdd$ and the settling time.
\end{problem}
\begin{solution}
Use the analytical approach with $\zeta=0.5$: the condition $1.5=2\cos\bigl(\wn T\sqrt{1-\zeta^{2}}\bigr)\ee^{-\zeta\wn T}$
is solved numerically, giving $\wn=4.32\,\mathrm{rad/s}$. Then $K=\ee^{-2\zeta\wn T}-0.5=\ee^{-0.432}-0.5=0.149$.
Check: $z^{2}-1.5z+0.649$ has the roots $0.75\pm0.294\,\jj$ with $|z|=0.806$; these give
$\zeta\wn=-\ln(0.806)/0.1=2.16$, $\wdd=\angle z/T=3.74\,\mathrm{rad/s}$, so $\wn=4.32$ and $\zeta=0.5$ as required.
The settling time is $T_s=4/(\zeta\wn)=1.85\,\mathrm{s}$. Note that the closed-loop poles always lie on the
vertical line $\Real[z]=0.75$ (the locus), so a single gain fixes only \emph{one} of the two quantities
($\zeta$ or $\wn$), never both.
\end{solution}

\begin{problem}[Overshoot of the three designs]
Using $PO=\ee^{-\pi\zeta/\sqrt{1-\zeta^{2}}}\times100$ estimate the overshoot and the settling time of designs
a)--c), and compare with the step response of the closed loop.
\end{problem}
\begin{solution}
\begin{center}
\renewcommand{\arraystretch}{1.25}
\begin{tabular}{c c c c c c}
Design & $K$ & $\zeta$ & $PO$ (formula) & $PO$ (simulated) & $T_s=4/\zeta\wn$\\\hline
a) & 0.23 & 0.30 & 37.2\% & 38.2\% & 2.5\,s\\
b) & 0.17 & 0.44 & 21.8\% & 22.3\% & 2.0\,s\\
c) & 0.10 & 0.70 & 4.6\% & 4.2\% & 1.6\,s
\end{tabular}
\end{center}
The closed loop has a pole at $z=1$ in $L(z)$, so the steady-state error to a step is zero in all three
cases. The formula and the sampled step responses agree to about one percentage point (the continuous-time
formula is only an approximation for a sampled system).
\end{solution}
